\begin{abstract}
% claim: AB01
Does the evidence in a motion observer's measurement correction improve anomaly detection beyond its prospective residual? We examine this question on one original CASPER right-arm recording containing eight scheduled speed-intervention intervals. An empirical categorical phase observer predicts twelve joint-position and joint-velocity measurements from a fixed periodic reference. The primary comparison adds a training-scaled Jensen--Shannon correction discrepancy while preserving the same observer state trajectory. Eleven deterministic score and control variants are fitted or specified using normal data, frozen before threshold-only calibration, and evaluated through continuous protected replay. All methods produce 897,119 finite eligible scores. Adding correction evidence increases unadjusted F1 from \RFone\ to \RCFone\ and slightly reduces normal alarm burden. Nevertheless, positive-point recall is only \RCRecallPct\%, with \RCEpisodeRate\ normal alarm episodes per eligible hour. Fresh alarms occur in all eight intervention intervals, showing why event hits must be interpreted alongside sparse point coverage and frequent normal alarms. An ordinary autoregression achieves higher F1 and average precision with more alarm episodes; a fixed periodic template attains the largest F1 with higher point false alarms. The result is a small within-recording correction-score effect, not evidence of a generally effective detector. Both primary variants retain progression, and no validated physical digital twin or cross-robot generalization is established. The accompanying local package preserves the complete comparison, source identities and deterministic table/figure generation; broader confirmation requires fresh independent evidence.
\end{abstract}
